


{\scriptsize\raggedright
\begin{longtable}{p{0.237\linewidth}p{0.237\linewidth}p{0.237\linewidth}p{0.237\linewidth}}
\toprule
No. & Where & Deviation & Effect \\
\midrule
\endhead
1 & Round 7, base & Round 7 branched from the head of \texttt{final2}, not from \texttt{submit-v3.1}: that tag does not exist (round 7b checked: no tag, no commit of the close-out prompt). \texttt{submit-v3} differs from the head of \texttt{final2} by two documents & none on a result; the close-out items that were missing were applied in round 7b (G0) \\
2 & Round 7, A bounded check & The enumeration used bounds that differ from the V18 text: allowance 16 to 40 units, a window of 30 steps, at most 2 earlier events (V18 text: 24 to 60 units, at most 3 earlier events). An earlier run with a window of 10 steps was too short for earlier events to share a window, so the unconditional form could not fail; it was replaced before anything was reported & run 1 (window 10): no result reported. Run 2 (reported in round 7): 1,139,264 configurations per policy, Invariant 1 and the conditional Invariant 2 hold, the one-event reserve fails the unconditional form only with 2 earlier events. Round 7b (V19) adds a third enumeration with allowances 40, 44, 48, a window of 70 steps and up to 5 earlier events \\
3 & Round 7, A1.5 & The ns-3 ACK-drop hook was cut: the ns-3 policy files are frozen, so the attack was simulated with the event engine \texttt{r7\_engine} (checked draw for draw against \texttt{r3\_event} and \texttt{micro\_event}), not in ns-3 & the A1 results are SIMULATED by the engine; ns-3 is NOT\_RUN for this item \\
4 & Round 7, documents & \texttt{scripts/doc\_style\_check.py} skips \texttt{PREREGISTRATION\_V18.md} because the file is frozen at its committed SHA-256 (its wording is not edited after the runs); V19 is not skipped and passes the check & the style tool has not checked the wording of V18 (it may contain words that the rules ban; for example, V18 uses the banned word once for a transmission) \\
5 & Round 7, B1 & The tag-coverage defect was found by the C benchmark, not by a model: at 2-byte pseudonyms two window entries of one vehicle can share a pseudonym, and a tag that does not cover the counter accepts the earlier entry (199,911 of 200,000 frames resolved). The tag now covers boot and counter as unsent associated data; the models were rerun & the fix is confined to the pseudonym variant (FRAME\_LAYOUT.md section 10.3); the ProVerif and Tamarin results quoted are for the fixed design \\
6 & Round 7b, G0 & The brief's "about 0.7 of a spreading-factor step" belongs to the 0 dBi case (1.85 dB). For the dipole case (4.00 dB) the same arithmetic gives 1.4 steps at the simulators' 2.8 dB per step; the paper uses the computed value & wording and key \texttt{r7b.g0.power.extra\_margin\_sf\_steps} \\
7 & Round 7b, V19 & Control F4 as written (the k = 1 policy violates the unconditional Invariant 2 with at most 3 earlier events) is not met: in the 70-step window the one-event reserve fails from 4 earlier events. The control works at the bound of 4 (k = 1 fails at an allowance of 40, k = 4 does not). The k = 4 policy fails from 5 earlier events (V19 expected 4) & reported as found in \texttt{ROUND7B\_RESULTS.md}; no check was changed after the run \\
8 & Round 7b, G3 & The round-7 report called the injective no-framing lemma "not proved". The exact status is a ProVerif trace of two base threads racing on the seen-set, plus the absence of any Tamarin quarantine model; a Tamarin model with an atomic seen-set was added and verifies both forms & the paper says the property holds with an atomic seen-set \\
9 & Round 7b, V19 & The size of the reserve (four events) was chosen from round-7 data; the registered confirmation uses fresh seeds and unchanged thresholds & confirmation, not a blind test; stated in V19 section 0 and in the paper \\
\bottomrule
\end{longtable}}

